%% introduction_draft.tex
%% Standalone draft — not yet integrated into main paper.
%% Flag to Neil before merging: check cross-references and bib keys.

\section{Introduction}
\label{sec:intro}

A foundational assumption in capital markets research is that investors interpret
public disclosures in a broadly consistent manner. Audit reports, in particular,
are designed to convey a standardized, institution-backed signal about the
reliability of financial statements. Yet the premise of homogeneous interpretation
sits uneasily alongside a growing body of evidence that political polarization
fundamentally reshapes how individuals process information, form beliefs, and
respond to institutions \citep{kempf2024, bordalo2023}. This paper asks whether
political polarization affects how investors interpret auditor-related information
--- and finds that it does, in economically meaningful ways.

We study auditor changes disclosed in Form 8-K (Item 4.01). These events are
well-suited to our question for three reasons. First, they are public, standardized
disclosures that all investors observe simultaneously, so differences in market
reactions must arise from differences in interpretation rather than information
access. Second, the signal is inherently ambiguous: a firm can change auditors for
benign reasons (e.g., audit committee rotation policy) or for troubling ones (e.g.,
disagreements over accounting estimates), and investors must form their own
assessments. Third, auditing is a regulated, institution-dependent activity, making
trust in regulatory infrastructure --- which is precisely what polarization erodes
--- central to how investors evaluate the disclosure.

Our theoretical framework builds on the heterogeneous-priors model of
\citet{kandel1995}. Investors observe a public audit signal but differ in their
prior beliefs due to differential trust in auditors and the institutions that
oversee them. Political polarization increases the dispersion of these priors,
which generates greater dispersion in posterior beliefs following an identical
signal. Under standard models of belief-driven trading \citep{harris1993},
greater posterior dispersion translates into more trading volume and larger
absolute price movements. The mechanism is distinct from information asymmetry:
polarization does not change what is disclosed, only how it is processed.

Empirically, we measure political polarization at the congressional district level
using the \citet{esteban1994} polarization index, applied to U.S. House election
returns from 2000 to 2023. We match firms to their headquarters state and construct
a state-by-year polarization measure that varies across both the cross-section and
time. Our main specification regresses absolute cumulative abnormal returns (CARs)
and abnormal trading volume over the $[-1, +1]$ window around auditor change
disclosures on the interaction of an event indicator with the polarization measure,
controlling for firm characteristics and including firm and year fixed effects.

We find that [PLACEHOLDER: insert main result once estimates are available ---
e.g., ``a one-standard-deviation increase in polarization is associated with
a X\% increase in absolute CAR and a Y\% increase in abnormal volume around
auditor change events.''] These effects are concentrated in \ldots [ambiguous
events / dismissals without disclosed reasons / Big-4 departures], consistent
with the model's prediction that polarization effects are amplified when signals
carry less inherent precision.

This paper makes three contributions. First, we contribute to the auditing and
capital markets literature by documenting that the market response to audit signals
depends on the political environment in which investors are embedded. Prior work
establishes that audit reports and auditor-related events convey information to
markets \citep{firth1978, dodd1984, defond2014}; we show that the same disclosure
generates systematically different reactions depending on investors' political
priors. Second, we contribute to the emerging literature on political polarization
in finance \citep{kempf2024}. Existing work focuses primarily on asset allocation,
partisan trading behavior, and analyst forecasts; we extend this literature to the
auditing context and to disagreement as a market outcome. Third, we contribute to
the theoretical literature on heterogeneous beliefs and trading
\citep{harris1993, kandel1995, hong2007} by providing an empirical setting where
the source of belief heterogeneity --- political polarization --- is plausibly
exogenous to firm-level reporting decisions and varies in well-defined ways over
time.

Our findings have implications for audit standard-setters and regulators. A
common motivation for expanded audit reporting requirements (e.g., critical audit
matters under PCAOB AS 3101) is that richer disclosures will reduce investor
uncertainty and improve market efficiency. Our results suggest a limit to this
logic: when investors differ in their priors about institutional credibility, more
information does not eliminate disagreement --- it may instead amplify it. The
effectiveness of disclosure-based regulation therefore depends on the political
and institutional context in which investors receive it.

The remainder of the paper proceeds as follows. Section~\ref{sec:lit} reviews
the related literature. Section~\ref{sec:theory} develops the theoretical
framework. Section~\ref{sec:design} describes the empirical design, data, and
variable construction. Section~\ref{sec:results} presents the results.
Section~\ref{sec:conclusion} concludes.
